\documentclass{article}

\usepackage{amsmath,amssymb}
\usepackage{xurl}
\usepackage[hidelinks]{hyperref}
\setlength{\emergencystretch}{2em}

% Shared commands for notes in docs/paper-gaps/.

\DeclareUnicodeCharacter{2080}{\ifmmode{}_{0}\else\textsubscript{0}\fi}
\DeclareUnicodeCharacter{2081}{\ifmmode{}_{1}\else\textsubscript{1}\fi}
\DeclareUnicodeCharacter{2082}{\ifmmode{}_{2}\else\textsubscript{2}\fi}
\DeclareUnicodeCharacter{2099}{\ifmmode{}_{n}\else\textsubscript{n}\fi}

\newcommand{\Lean}{\textsc{Lean}}
\ExplSyntaxOn
\NewDocumentCommand{\leanid}{m}{
  \ifmmode
    \textsf{\detokenize{#1}}
  \else
    \group_begin:
    \urlstyle{sf}
    \tl_set:Nn \l_tmpa_tl {#1}
    \tl_replace_all:Nnn \l_tmpa_tl {\_} {_}
    \exp_args:NV \path \l_tmpa_tl
    \group_end:
  \fi
}
\ExplSyntaxOff
\providecommand{\tr}{\operatorname{tr}}
\newcommand{\tnleanIssueBase}{https://github.com/LionSR/TNLean/issues/}
\newcommand{\tnleanPrBase}{https://github.com/LionSR/TNLean/pull/}
\newcommand{\tnleanDocBase}{https://sirui-lu.com/TNLean/docs/}
\newcommand{\ghissue}[1]{\href{\tnleanIssueBase#1}{GitHub issue \##1}}
\newcommand{\ghpr}[1]{\href{\tnleanPrBase#1}{PR \##1}}
\newcommand{\doclink}[1]{\href{\tnleanDocBase#1.html}{\path{#1}}}

% Dirac notation and the identity operator, as in blueprint/src/macros/common.tex.
% \providecommand keeps note-local definitions authoritative.
\usepackage{dsfont}
\providecommand{\ket}[1]{|#1\rangle}
\providecommand{\bra}[1]{\langle #1|}
\providecommand{\braket}[2]{\langle #1 | #2 \rangle}
\providecommand{\ketbra}[2]{|#1\rangle\!\langle #2|}
\providecommand{\Id}{\mathds{1}}
\providecommand{\one}{\mathds{1}}
\providecommand{\C}{\mathbb{C}}
\providecommand{\MN}[1]{M_{#1}(\C)}
\providecommand{\MD}{M_D(\C)}

% External referencing for paper sources.  The build helper writes sanitized
% auxiliary files under build/paper-aux/ and excludes duplicated source labels.
\usepackage{xr}

% Import a generated paper auxiliary file with a caller-supplied label prefix.
% The candidate paths support builds from docs/paper-gaps/, the repository root,
% and build/paper-aux/ itself.
\newcommand{\importpaperaux}[2]{%
  \IfFileExists{../../build/paper-aux/#2.aux}{%
    \externaldocument[#1]{../../build/paper-aux/#2}%
  }{%
    \IfFileExists{build/paper-aux/#2.aux}{%
      \externaldocument[#1]{build/paper-aux/#2}%
    }{%
      \IfFileExists{#2.aux}{%
        \externaldocument[#1]{#2}%
      }{}%
    }%
  }%
}

% General external reference macros
\newcommand{\paperref}[2]{\ref{#1-#2}}
\newcommand{\papereqref}[2]{(\ref{#1-#2})}

% CPSV16 source (1606.00608 / MPDO-22-12-17-2.tex) external document and shortcuts
\importpaperaux{cpsv16-}{1606.00608/MPDO-22-12-17-2-tnlean}
\newcommand{\cpsvref}[1]{\paperref{cpsv16}{#1}}
\newcommand{\cpsveqref}[1]{\papereqref{cpsv16}{#1}}

% Verdict marker: \gapnote{<kind>}{<status>}. Kind is the policy
% classification (clarification, local-correction, scope-restriction,
% unfaithful, false-source, open-gap); status is open, wip (elimination
% underway), resolved, or historical. Machine-read by the site index;
% prints a verdict line.
\newcommand{\gapnote}[2]{\par\noindent\textbf{Verdict:} #1 (#2).\par}


\title{The Positive-Gap Boundary for Distance- and Length-Independent MPDO
Correlations}
\author{}
\date{2026-08-29}

\begin{document}
\maketitle
\gapnote{scope-restriction}{resolved}

\section{Source formula and mixed-state reading}

Equation~\cpsveqref{Corr} of
Ref.~\cite{Cirac2016MPDO_arXiv} writes the two-observable contraction of a
periodic matrix product vector as
\begin{align}
    \langle V|O_1O_2|V\rangle
        &=\tr\!\left(
          \mathbb E_{O_1}\mathbb E^{\,n_2-n_1-1}
          \mathbb E_{O_2}\mathbb E^{\,N+n_1-n_2-1}\right).
    \label{eq:cpsv16_mpdo_zcl_correlation_length_boundary_source_correlation}
\end{align}
Here $N$ is the ring length and $n_1<n_2$ are the marked sites.
Definition~\cpsvref{DefinitionZCL} gives the mixed-state zero-correlation-length
condition, and the following sentence states that the corresponding MPDO
correlation functions are independent of length
\cite[lines~735--742]{Cirac2016MPDO_arXiv}.

For an MPO tensor $M$, closing the physical indices gives
\begin{align}
    \mathcal T_M
        &=\sum_i M^{ii},
    \label{eq:cpsv16_mpdo_zcl_correlation_length_boundary_physical_trace_transfer}
\end{align}
and inserting a one-site observable $O$ in the same contraction gives
\begin{align}
    \mathbb E_O
        &=\sum_{i,j}O_{ji}M^{ij}.
    \label{eq:cpsv16_mpdo_zcl_correlation_length_boundary_observable_insertion}
\end{align}
Thus the mixed-state counterpart of
Eq.~\ref{eq:cpsv16_mpdo_zcl_correlation_length_boundary_source_correlation}
is
\begin{align}
    C_M(O_1,O_2;g_{\mathrm{mid}},g_{\mathrm{wrap}})
        &=\tr\!\left(
          \mathbb E_{O_1}\mathcal T_M^{g_{\mathrm{mid}}}
          \mathbb E_{O_2}\mathcal T_M^{g_{\mathrm{wrap}}}\right),
    \label{eq:cpsv16_mpdo_zcl_correlation_length_boundary_mixed_correlation}
\end{align}
where
\begin{align}
    g_{\mathrm{mid}}
        &=n_2-n_1-1,
    \label{eq:cpsv16_mpdo_zcl_correlation_length_boundary_middle_gap}\\
    g_{\mathrm{wrap}}
        &=N+n_1-n_2-1.
    \label{eq:cpsv16_mpdo_zcl_correlation_length_boundary_wrapping_gap}
\end{align}
Equivalently, the ring length is
\begin{align}
    N
        &=g_{\mathrm{mid}}+g_{\mathrm{wrap}}+2.
    \label{eq:cpsv16_mpdo_zcl_correlation_length_boundary_ring_length}
\end{align}

The contraction in
Eq.~\ref{eq:cpsv16_mpdo_zcl_correlation_length_boundary_mixed_correlation}
is unnormalized. The source introduces normalization of finite-chain MPDOs
only later, at lines~792--793, for entropic quantities. The unnormalized
reading therefore requires no nonvanishing hypothesis absent from
Definition~\cpsvref{DefinitionZCL}.

\section{Positive varied gaps}

Definition~\cpsvref{DefinitionZCL} is the literal idempotence equation
\begin{align}
    \mathcal T_M^2
        &=\mathcal T_M.
    \label{eq:cpsv16_mpdo_zcl_correlation_length_boundary_idempotence}
\end{align}
It implies
\begin{align}
    \mathcal T_M^r
        &=\mathcal T_M,
        \qquad r\geq 1.
    \label{eq:cpsv16_mpdo_zcl_correlation_length_boundary_positive_power}
\end{align}
For the length-independence assertion, fix an arbitrary middle gap
$g_{\mathrm{mid}}\geq0$ and compare two positive wrapping gaps
$g_{\mathrm{wrap}},g'_{\mathrm{wrap}}>0$. Substitution of
Eq.~\ref{eq:cpsv16_mpdo_zcl_correlation_length_boundary_positive_power} for
the two wrapping powers gives
\begin{align}
    C_M(O_1,O_2;g_{\mathrm{mid}},g_{\mathrm{wrap}})
        &=C_M(O_1,O_2;g_{\mathrm{mid}},g'_{\mathrm{wrap}}).
    \label{eq:cpsv16_mpdo_zcl_correlation_length_boundary_positive_wrap_gap_independence}
\end{align}
The middle power is identical on both sides and need not have a positive
exponent. By
Eq.~\ref{eq:cpsv16_mpdo_zcl_correlation_length_boundary_ring_length}, varying
the wrapping gap at fixed middle gap varies the ring length while keeping the
observable separation fixed. Thus
Eq.~\ref{eq:cpsv16_mpdo_zcl_correlation_length_boundary_positive_wrap_gap_independence}
is the positive-wrapping-gap scope restriction of the length-independence
claim for MPDO correlations at lines~740--742. It includes adjacent marked
sites, for which $g_{\mathrm{mid}}=0$.

For the distance comparison, take two pairs of positive complementary gaps
$g_{\mathrm{mid}},g_{\mathrm{wrap}}>0$ and
$g'_{\mathrm{mid}},g'_{\mathrm{wrap}}>0$. Applying
Eq.~\ref{eq:cpsv16_mpdo_zcl_correlation_length_boundary_positive_power} to
all four powers gives
\begin{align}
    C_M(O_1,O_2;g_{\mathrm{mid}},g_{\mathrm{wrap}})
        &=C_M(O_1,O_2;g'_{\mathrm{mid}},g'_{\mathrm{wrap}}).
    \label{eq:cpsv16_mpdo_zcl_correlation_length_boundary_positive_gap_independence}
\end{align}
If the two gap sums are equal, the ring length in
Eq.~\ref{eq:cpsv16_mpdo_zcl_correlation_length_boundary_ring_length} is fixed
and
Eq.~\ref{eq:cpsv16_mpdo_zcl_correlation_length_boundary_positive_gap_independence}
gives correlation independence under changing the observable distance while
the observables are nonadjacent on both complementary arcs. This is the
positive-complementary-gap scope restriction of the implication from
equation~\cpsveqref{Corr} to CID in the discussion at lines 433--500 of
Ref.~\cite{Cirac2016MPDO_arXiv}.

Positivity is required only for an exponent that is replaced using
Eq.~\ref{eq:cpsv16_mpdo_zcl_correlation_length_boundary_positive_power}. At a
zero gap,
\begin{align}
    \mathcal T_M^0
        &=\Id,
    \label{eq:cpsv16_mpdo_zcl_correlation_length_boundary_zero_power}
\end{align}
and
Eq.~\ref{eq:cpsv16_mpdo_zcl_correlation_length_boundary_idempotence} does not
identify $\Id$ with $\mathcal T_M$. Hence the length theorem permits an
unchanged zero middle gap but does not compare a zero wrapping gap with a
positive one. The distance theorem, which varies both gaps, requires all four
compared gaps to be positive. The unrestricted source/CID reading involving a
zero varied complementary gap is not claimed.

\section{Formal status}

The observable insertion in
Eq.~\ref{eq:cpsv16_mpdo_zcl_correlation_length_boundary_observable_insertion}
is \leanid{MPOTensor.verticalLoopWith}. The correlation in
Eq.~\ref{eq:cpsv16_mpdo_zcl_correlation_length_boundary_mixed_correlation}
is \leanid{MPOTensor.periodicTwoPointCorrelation}, and
Eq.~\ref{eq:cpsv16_mpdo_zcl_correlation_length_boundary_positive_wrap_gap_independence}
is
\leanid{MPOTensor.periodicTwoPointCorrelation_positiveWrapGaps_independent}.
This length-independence theorem assumes only literal physical-trace
idempotence and positivity of the two compared wrapping gaps; the fixed middle
gap may be zero. It introduces no positivity, normalization, nonvanishing, or
equal-total-length condition on the MPO tensor.

Eq.~\ref{eq:cpsv16_mpdo_zcl_correlation_length_boundary_positive_gap_independence}
is \leanid{MPOTensor.periodicTwoPointCorrelation_positiveGaps_independent}.
This fixed-ring distance corollary requires positivity of the four compared
complementary gaps; equal gap sums give the nonadjacent-on-both-arcs CID
specialization.

The scope restriction is resolved by requiring positivity precisely for each
gap exponent that is varied. No invariance replacing a zeroth power by a
positive power is claimed.

\bibliographystyle{alpha}
\bibliography{references}

\end{document}
